\section{Introduction}

The operational carbon associated with training depends on electricity
consumption and when that electricity is used~\cite{patterson2021carbon}.
Carbon-aware systems exploit temporal, spatial, and allocation
flexibility~\cite{hanafy2023carbonscaler,sukprasert2024limitations}. On a shared
batch cluster, however, a user controls a \emph{request}, while Slurm determines
admission, placement, and start time. A resource choice may take effect hours
later, when the carbon intensity differs from that observed at submission.

Node count couples three consequences. Fewer nodes can reduce node-hours under
sublinear strong scaling, but lengthen execution and increase the number of
queue admissions needed by a long run. A smaller request may start sooner,
yet overlap a less favorable carbon interval. Near completion, a faster final
chunk may avoid another queue wait. The useful decision therefore depends on
remaining work, available time, and the execution window of each request.

The central question is \emph{when adapting requests provides value beyond a
good fixed scale}. With constant carbon intensity, no waiting, and no chunk
overhead, we show that mixed-scale execution lies in the convex hull of
fixed-scale time--carbon points. Scaling inefficiency alone offers no gain
below that frontier. Under deadlines, a preplanned within-run mix can still
outperform a per-run randomized fixed policy without new queue observations.
We therefore retain both fixed mixtures and a plan chosen once at submission
when testing the value of online adaptation.

\name{} exposes a completion-budget slider and makes repeated node-count decisions without
modifying Slurm. Its main actor--critic policy uses visible queue history,
remaining work and time, action execution descriptors, and a causal CI forecast
sequence. It requires no explicit wait predictor: continuation costs are
learned from realized replay outcomes, and each checkpoint reveals how much
budget remains. Hidden admissions still limit predictability and any deadline
claim. We compare against explicit rollout planning and a matched learned policy
that commits its full sequence before submission, separating dynamic feedback
from the benefit of a preplanned mix.

Complete-node power is unavailable, so we separate execution evidence from
carbon estimates. The power model has three factors: a normalization
reference, an unknown workload coefficient, and a scale coefficient. For a fixed
execution log, carbon is affine in the scale-model parameter. Two endpoint
comparisons thus characterize robustness within the declared interval,
without determining absolute power or repeating scheduler simulations.
We use this property to train against both endpoint objectives and report
where comparative conclusions change.

Our evaluation uses published AMSP performance profiles and two historical
arrival streams to define controlled contention scenarios. Hardware labels do
not determine admission once resource demands and scheduling rules are fixed;
this abstraction does not predict how old jobs would run on new hardware.

The paper makes three contributions:
\begin{itemize}
\item A completion-budget formulation and fixed-scale null case that separate
scaling efficiency, preplanned mixing, and online adaptive value.
\item A user-level controller combining realized budget feedback, public queue
history, and undiscounted constrained learning without an admission predictor.
\item An exposure-based sensitivity analysis and a matched evaluation against
fixed, plan-once, and replanning baselines.
\end{itemize}
